\bookchapter{Appendice A: Flow “tutto-per-tutto” per lo sballaggio del nuovo telefono}

\booksection{“Trilogia” di verifica allo sballaggio:}

\begin{itemize}
  \item \textbf{Aspetto}: controlla corpo e schermo per graffi, ammaccature.
  \item \textbf{Dati}: verifica che IMEI/numero di serie sulla scatola corrisponda a quello della fattura.
  \item \textbf{Funzioni}: fai una chiamata per provare auricolare e microfono, scatta una foto per controllare la fotocamera.
\end{itemize}

\booksection{“Promemoria chiave” per attivazione e configurazione iniziale:}

\begin{itemize}
  \item \textbf{Rete}: prima connettiti a un Wi-Fi stabile.
  \item \textbf{Account}: accedi al tuo Apple ID, account Huawei, account Xiaomi, ecc.
  \item \textbf{Impostazioni di base}: attiva "Trova il mio telefono", imposta una password di blocco schermo robusta.
  \item \textbf{eSIM}: se il tuo telefono supporta eSIM (come iPhone, alcuni flagship Android), puoi aggiungere un piano eSIM nelle impostazioni, senza inserire una SIM fisica. All'estero l'eSIM è particolarmente comoda, puoi comprare direttamente un piano eSIM di un operatore locale.
  \item \textbf{Migrazione dati}: quando trasferisci i dati dal vecchio telefono, assicurati che entrambi abbiano carica sufficiente. L'iPhone usa "Quick Start", Android usa l'app di migrazione ufficiale.
\end{itemize}

\booksection{Installazione app “sicurezza prima”:}

\begin{itemize}
  \item \textbf{Prima scelta: store ufficiale}: Apple usa App Store, Android usa lo store app integrato nel telefono.
  \item \textbf{Non dare permessi a caso}: quando l'installazione chiede permessi, chiediti "quest'app ne ha davvero bisogno?"
\end{itemize}
